\section{Discussion and Methodological Guidelines}
\label{sec:discussion}

The results of our methodological audit have significant implications for both researchers and practitioners in dynamic model merging. Based on our empirical revelations, we outline a set of constructive guidelines to improve evaluation standards and guide deployment decisions.

\subsection{Guidelines for Rigorous Evaluation (For Researchers)}
To prevent future false progress and ensure scientific clarity, we advocate for the following practices in the model merging literature:
\begin{enumerate}
    \item \textbf{Proper Baseline Regularization and Initialization:} When introducing a new dynamic merging or routing method, authors must compare against parametric baselines that are properly initialized (such as maximum-entropy zero-initialization) and swept across multiple regularization strengths ($\lambda$). Comparing against an unregularized, randomly-initialized baseline on small datasets is a straw-man comparison.
    \item \textbf{Evaluate Across Dual Data Regimes:} A method's performance must be contextualized by the size of the calibration set ($N_{\text{cal}}$). Training-free geometric projection techniques should be categorized as sample-efficient priors, while parametric routers should be evaluated in high-data regimes to establish their performance upper bounds.
    \item \textbf{Track Intermediate Representation Quality:} Evaluation should not be restricted to end-to-end task accuracy. Tracking intermediate representations (such as layer-wise cosine similarity to prototypes) provides crucial mechanistic insights, exposing phenomena like representational lag and feature warping that are masked by final logits.
\end{enumerate}

\subsection{Deployment Decision Matrix (For Practitioners)}
For practitioners deploying multi-task systems on edge devices, the choice of routing architecture depends heavily on the available data budget, computational constraints, and environmental noise levels:

\paragraph{Ultra-Low Calibration Budgets ($N_{\text{cal}} \le 128$).}
When data collection is extremely constrained, training-free centroid-based methods like SABLE act as highly robust geometric priors, completely bypassing the risk of parametric overfitting and eliminating optimization overhead.

\paragraph{Sufficient Calibration Budgets ($N_{\text{cal}} \ge 1000$).}
When calibration data is abundant, a simple stateless classical linear router is mathematically superior, computationally lighter, and easier to implement. Under this regime, the unregularized router achieves competitive accuracy and executes sharp, instantaneous ensembling decisions without serving-time ODE overhead.

\paragraph{Stateful Kinetics for Feedback Stabilization.}
Stateful, continuous-time kinetics (e.g., ChemMerge) should be preferred when serving in deep networks under heavy, unbalanced activation noise. Because our high-fidelity sandbox models SVHN (Task 3) under an elevated noise level ($\sigma_3 = 1.20$), the closed-loop stabilizing dynamics of ChemMerge act as a beneficial temporal low-pass filter (closed-loop stateful inertia). This closed-loop feedback dynamically corrects early routing decisions and prevents premature commitment to a wrong task manifold, achieving a superior performance ceiling. However, if the downstream environment features a highly balanced task noise profile (e.g., $\sigma_k = 0.05$ or $0.15$ across all tasks), the need for temporal low-pass filtering is virtually eliminated. Under such balanced regimes, the stabilization premium of ChemMerge narrows to statistical insignificance, making the simpler, stateless classical router the globally optimal deployment choice. Furthermore, in non-stationary streaming applications characterized by sudden task transitions, this stateful inertia introduces a representational lag (hysteresis) that can lead to transient classification errors immediately after a task switch. In such highly non-stationary environments, we recommend implementing an activation-based ``state reset'' trigger to flush the ODE state and restore routing agility. Mathematically, this can be integrated into the continuous reaction kinetics (Equation 1) by introducing a gating variable $g(t) = \sigma(\gamma (\|\nabla h(t)\| - \theta))$, which detects sharp representational transitions (where $\|\nabla h(t)\|$ is the activation derivative across layers or adjacent samples, $\theta$ is a threshold, and $\gamma$ is a scaling factor). When a transition is detected ($g(t) \approx 1$), the concentration dynamics are overridden:
\begin{equation}
\frac{d C_k(t)}{dt} = - K_{\text{decay}} C_k(t) + R_k(t) + \delta \cdot g(t) \left( C_{0} - C_k(t) \right)
\end{equation}
where $C_0$ is the uniform concentration prior, and $\delta \gg 1$ is a large penalty that rapidly resets the system's concentrations to their maximum-entropy state, effectively flushing the stateful memory while preserving the closed-loop stabilization premium during stable phases.

\paragraph{Discussion of Sandbox Noise Unbalance and Expert Realism.}
We must explicitly note that the extreme asymmetric noise profile of our sandbox (where the SVHN expert achieves only $22.80\%$ standalone accuracy under a massive noise variance $\sigma_3 = 1.20$) creates an artificially difficult environment that heavily penalizes open-loop parametric routers and strongly favors the stateful, continuous ensembling of ChemMerge. In realistic multi-task serving deployments, engineers rarely deploy expert adapters with such poor, near-guessing standalone performance. Under a more realistic, balanced noise profile where all expert adapters achieve robust, stable standalone performance ($>80\%$), the premium for closed-loop ensembling weight trajectory smoothing would narrow significantly, leaving the simpler, stateless classical linear router as the globally optimal, latency-efficient serving choice across all data regimes.

\paragraph{Quantitative Serving-Time Complexity Analysis.}
To ground this decision matrix in system-level hardware performance and serving-time efficiency, we present a comparative quantitative complexity analysis in Table~\ref{tab:serving_complexity}. We map the serving-time parameter complexity, FLOPs per inference, evaluation schedule, and sequential layer-wise overhead across all evaluated gating architectures. Here, $L' = 11$ is the number of dynamic blending layers (Layers 4 to 14) and $N_{\text{step}}$ is the number of discretized Euler steps (with $N_{\text{step}} \ge 1$) per layer.

\begin{table*}[ht]
\caption{Quantitative Complexity Analysis of Serving-Time Routers}
\label{tab:serving_complexity}
\centering
\begin{footnotesize}
\setlength{\tabcolsep}{3.0pt}
\begin{tabular}{lcccc}
\toprule
\textbf{Method} & \textbf{Parameters} & \textbf{FLOPs per Inference} & \textbf{Gating Evaluation} & \textbf{Sequential Overhead} \\
\midrule
Stateless Invariant & $K \cdot D$ ($768$) & $O(K \cdot D)$ & Once (at Layer 3) & None \\
Layer-Wise Classical & $L' \cdot K \cdot D$ ($8,448$) & $O(L' \cdot K \cdot D)$ & Layer-wise & Low (Linear proj.) \\
Stateless SABLE & $0$ ($K \cdot D$ centroids) & $O(L' \cdot K \cdot D)$ & Layer-wise & Medium (Cosine sim.) \\
Stateful ChemMerge & $0$ ($K \cdot D$ centroids) & $O(L' \cdot (K \cdot D + K \cdot N_{\text{step}}))$ & Layer-wise & High (ODE + clamping) \\
\bottomrule
\end{tabular}
\end{footnotesize}
\end{table*}

\subsection{Limitations, Trade-offs, and Future Work}
\label{subsec:limitations_future_work}

While our methodological audit provides crucial conceptual and mechanistic insights into activation-space model merging, we must transparently outline several limitations, trade-offs, and avenues for future research:

\begin{enumerate}
    \item \textbf{Layer-Invariant vs. Layer-Wise Parametric Routing Trade-off:} SABLE and ChemMerge are true \emph{layer-wise dynamic routers} that compute independent ensembling coefficients at every layer. In Section~\ref{subsec:layerwise_ablation}, we executed a thorough ablation by implementing a true \emph{layer-wise classical router} with separate gating heads at each layer, scaling the parameter count from $768$ to $8,448$ parameters. 
    
    Our results demonstrate that the layer-wise classical router achieves highly smooth routing trajectories (Jitter as low as $0.0068$), proving that parametric routers do not inherently suffer from routing jitter. However, keeping the parametric gating layer-invariant acts as a powerful structural regularizer (inductive bias) that reduces parameter complexity by a factor of 11, making small-sample calibration extremely stable while maintaining high serving accuracy. Further exploring parameter-efficient layer-wise parametric routing (e.g., using shared low-rank gating weights across layers) represents a valuable direction for future research.
    
    \item \textbf{Absence of Train/Test Distribution Shift in Sandbox:} Our Analytical Coordinate Sandbox (ICS) assumes that the calibration (train) set and downstream test set are drawn from the same underlying distribution. Under this identically distributed setting and abundant calibration data ($N_{\text{cal}}=4000$), our unregularized Softmax Router achieves the absolute highest accuracy, while the heavily regularized router ($\lambda = 10^{-2}$) suffers from constraint bias. 
    
    In realistic edge deployments, however, the downstream test stream is almost always characterized by temporal drift, task-order non-stationarity, or heteroscedastic noise. In such non-stationary environments, an unregularized parametric router would severely overfit to the clean calibration split, suffering catastrophic generalization drops on the test stream. Under such realistic distribution shifts, proper regularization (or training-free priors) is mandatory to survive, making our proposed regularized baseline ($\lambda = 10^{-4}$) or stateful continuous-time priors even more critical.
    
    \item \textbf{Terminology Framing and Simplicity:} We acknowledge that our introduced formulations of \emph{Maximum-Entropy Zero-Initialization} and \emph{Proper L2 Regularized Calibration} are functionally equivalent to standard zero-initialization and L2 weight decay. We employ these high-level framings not to obscure simple concepts, but to emphasize their physical and mathematical properties: zero-initialization acts as a maximum-entropy prior that guarantees complete representational symmetry and unbiased starting manifolds, while L2 regularization actively restricts the search space under under-determined constraints. Framing these standard techniques through this lens highlights why they serve as the mathematically correct baselines that should never be omitted. Furthermore, this regularization perspective connects deeply with the classical sparse Mixture-of-Experts (MoE) literature, where specialized regularization techniques (such as auxiliary load-balancing losses, entropy regularization, or router noise injection \citep{shazeer2017outrageously}) are routinely employed to prevent router collapse and representational imbalance. Drawing this link shows that preventing routing collapse via proper regularization is a foundational, well-established principle in deep learning that was simply overlooked or omitted in the dynamic model merging sub-field.
    
    \item \textbf{Generalizability from Synthetic Sandbox to Large Foundation Models:} While our high-fidelity Analytical Coordinate Sandbox (ICS) is carefully parameterized to mirror the representation cones and singular value spectra of real-world vision and language models, simulations inherently simplify raw data geometry. Specifically, our Toeplitz covariance injection ($\Sigma_{i,j} = \rho^{|i-j|}$) models representation anisotropy elegantly, but real foundation model cones are highly non-symmetric. For future research, we explicitly suggest a concrete non-parametric covariance modeling technique: estimating empirical covariance matrices directly from raw pre-trained token embeddings across diverse corpora, and using the resulting eigenvectors to inject a more realistic, heteroscedastic noise-shaping prior.
    
    Furthermore, while we validated generalizability on a pre-trained \textbf{BERT-Tiny} model with custom LoRA adapters on actual GLUE datasets (SST-2 vs. QQP) and confirmed that classical parametric routers are highly effective ($62.50\%$ serving accuracy compared to SABLE's $60.00\%$) under a calibration budget of 500 samples, these classification tasks map to highly separated subspaces. For generative tasks (e.g., text summarization vs. code generation on large LLaMA or Mistral-based adapters), task embeddings exhibit much larger geometric overlap due to shared syntactic structures. In such generative environments, the routing boundary is non-trivial, meaning the small-sample overfitting bottleneck could re-emerge as a critical hazard. Under such dense overlaps, we hypothesize that the inductive geometric priors of SABLE and ChemMerge will maintain a high accuracy premium, whereas classical parametric routers will require stronger regularization or a hybrid hard-soft gating strategy. Evaluating these dynamics under generative, instruction-tuned workloads represents an exciting and essential direction for future deployment research.

    \item \textbf{Direct Classifier Logit Blending and Label Space Alignment Constraints:} In our real-world BERT-Tiny experiments, the joint multi-task serving model computes predictions by taking a weighted sum of the task-specific classifiers' logit outputs: $\text{logits} = \alpha_0 \cdot \text{classifier}_0(\text{pooled}) + \alpha_1 \cdot \text{classifier}_1(\text{pooled})$. While mathematically convenient for this binary classification setup, this design assumes that all task-specific classifiers output logits of the exact same dimensionality. If we were to serve tasks with mismatched label spaces (e.g., combining a 2-class sentiment analysis task with a 3-class MNLI subtask or a 10-class subtask), this direct blending equation would crash immediately due to a shape mismatch. In realistic multi-task serving environments, a general-purpose routing framework must instead route input samples directly to their respective task classifier heads without logit-level blending, or implement dynamic label space adapters to handle task-specific label spaces.
    
    \item \textbf{BERT-Tiny Scale and Expert Standalone Performance Bounds:} Due to local compute constraints, our foundation model validation was restricted to a compact pre-trained BERT-Tiny model (4 encoder layers, hidden size 128) evaluated on custom LoRA adapters trained in a highly data-constrained regime (standalone test accuracies of $58.80\%$ on SST-2 and $65.60\%$ on QQP). While this toy architecture serves as a valuable proof-of-concept, it does not fully capture the high-dimensional activation manifolds, representation geometry, or scale-dependent memory and compute bottlenecks of modern multi-billion parameter foundation models (such as LLaMA, Mistral, or ViT-B/16). Broadening our empirical findings to larger, state-of-the-art models with fully converged standalone experts represents a crucial avenue for future work to confirm if our methodological insights scale proportionally.

    \item \textbf{Designing True Closed-Loop Parametric Routers:} Throughout our analysis, we showed that the superior performance ceiling of ChemMerge ($76.90\%$ in the sandbox) is primarily due to its closed-loop feedback structure, which dynamically re-evaluates ensembling coefficients based on intermediate layer activations as they are cleaned up. Conversely, our classical parametric router is a feedforward, open-loop controller that makes a single routing decision based on early representation features. A promising future direction would be to combine the learning capacity of parametric models with the dynamic stabilization of closed-loop systems. Specifically, one could design a \emph{closed-loop parametric router} where the gating head receives intermediate layer representation feedback. Such a network could be trained end-to-end (e.g., using Recurrent Neural Networks, Transformer layers, or simple gating networks that ingest representations across different depths) to dynamically update routing weights $\boldsymbol{\alpha}^{(l)}$ as activations pass through the network, allowing the system to learn the optimal feedback controller directly from data rather than relying on hand-crafted continuous-time kinetics metaphors.

    \item \textbf{Evaluating on Fully Converged Experts:} In our real-world validation, the pre-trained BERT-Tiny model utilized custom LoRA experts that were under-fitted due to data constraints. Under-fitted experts might generate noisier activation manifolds, which could potentially distort ensembling behavior. Future extensions should evaluate routing architectures on fully converged, high-performing experts to ensure that routing accuracy is not fitting to spurious noise or representational artifacts, thereby confirming that the relative performance dynamics remain invariant to expert training budgets.

    \item \textbf{Asymmetry in Embedding-Level vs. Layer-Wise Gating:} In our BERT-Tiny experiments, the classical parametric router is evaluated as a stateless, embedding-level gating model (at Layer 0) for efficiency and simplicity, whereas SABLE and ChemMerge compute routing decisions dynamically layer-by-layer. This architectural asymmetry means our classical router operates in a feedforward, open-loop manner at the input level, while the training-free baselines operate layer-wise. Investigating layer-wise parametric gating in real foundation models—by training tiny routing heads at each layer—could bridge this structural gap and potentially recover the closed-loop stabilization premium observed in synthetic environments.
\end{enumerate}

\section{Conclusion}
\label{sec:conclusion}

In this work, we have presented a rigorous methodological audit and experimental deconstruction of dynamic model merging from the perspective of \emph{The Methodologist}. By applying Occam's razor, we have demonstrated that the widely reported failure of classical parametric routers is not an inherent representational limitation, but a confounding artifact of weak baseline tuning under small-sample constraints. 

Our proposed, properly regularized, maximum-entropy zero-initialized classical Softmax router achieves highly competitive performance in high-data regimes, successfully outperforming SABLE with significantly less computational complexity. While stateful continuous-time routing (ChemMerge) maintains a superior performance ceiling due to its closed-loop feedback structure, our audit reveals that stateless classical gating represents an incredibly efficient and simple alternative when calibration data is available. We hope these insights restore simplicity, clarity, and methodological rigor to the field of dynamic model merging.
